\chapter{A Difference Appears}

\section{The first fact}

Before a thing can be measured it must first be told apart from what it is not, and on this face that telling-apart is a fact
of geometry stripped to its bones. There is no metric yet, no length, no clock; there is a smooth manifold $M$ and, on it,
two points $p$ and $q$ with $p \neq q$. That inequality is the whole of the ground floor. It says nothing about how far
apart the points lie --- distance is a later structure, laid on top --- only that the smooth structure already tells them
apart, that a chart $\varphi : U \to \mathbb{R}^n$ assigns them different coordinates, $\varphi(p) \neq \varphi(q)$. The
first fact of the physical instrument is this bare nondegeneracy: a separation that exists before anything has been assigned
a size.

It is worth being exact about how thin this is, because the whole discipline of the book is the refusal to inflate it. A
manifold is, before all else, a Hausdorff space: distinct points possess disjoint neighborhoods. That axiom, and it alone,
is what ``distinguishable'' means here. Two points are told apart when some open set contains the one and not the other ---
a purely topological verdict, decided without a ruler, prior to any $g_{\mu\nu}$. We have not yet said that $p$ and $q$ are a
meter apart, or a nanometer, or that one lies to the left of the other. We have said only that the space does not identify
them: there is a neighborhood that separates them, and the existence of that neighborhood \emph{is} the difference. The
Hausdorff separation is the geometric form of a $\emph{this}$ that is not a $\emph{that}$.

Notice what has \emph{not} appeared, and how deliberate each absence is. There is no inner product, so there is no
$\lVert p - q \rVert$ --- indeed $p - q$ is not even defined, since points of a manifold do not subtract. There is no
orientation, so there is no sense in which one point precedes the other. There is no measure, so there is no volume, no
probability, no ``how much.'' The manifold at this tier carries exactly one relation, $\neq$, and that relation is symmetric
and blind: it reports two-ness and stops. Everything the physicist will later pour onto this space --- the metric, the
connection, the field --- is structure \emph{added} to the bare fact of distinctness, and none of it is present yet. The
first fact is the emptiest true thing the geometry can say.

And here the first of the book's two working verbs enters where it is native. To \emph{tange} is to select --- to pick out
one point against another by the characteristic that parts them. On this face a tange is a \emph{contravariant} act: it lays
a label on the world, a coordinate $x^\mu$, a chart value, our chosen index for \emph{this} point and not \emph{that} one.
The label transforms against the world, $x^\mu \mapsto x'^\mu = \frac{\partial x'^\mu}{\partial x^\nu}\,x^\nu$, exactly
because it is ours and not the world's; a change of frame reshuffles our labels while the points sit still. To tange $p$
from $q$ is to hold a contravariant index that discriminates them. The partner verb, to \emph{funge} --- to gather by
sameness, the covariant act that counts what the world does --- is not yet available, and cannot be, because you cannot bag
points by a shared characteristic until you have first told them apart. Tange precedes funge here as selection precedes
counting, and the order will hold the length of the book.

This lands the reader on the one fact the ground floor exists to establish: \emph{resolution is logically upstream of
measurement}. You must have two distinct points before you can have a distance between them; you must resolve a separation
before you can size it. The metric $g_{\mu\nu}$, when it comes, will act on pairs of tangent vectors that themselves presume
the points are already told apart. So the smooth structure's power to separate is not a small preliminary to the physics; it
is its precondition. A field cannot be read on a space that cannot first tell its points apart, any more than a scale can be
read whose marks blur into one. The mark must be a distinctness before it can be a magnitude.

\section{The symbol as a label}

The first fact was a separation; the natural next question is where the geometry \emph{writes it down}. Having told $p$ from
$q$, where does it record the distinction so as to find it again? The answer is the hinge of the chapter, and it is a small
shock: nowhere in particular, because the distinction is already its own record. A coordinate is not a fact about a point
stored alongside the point; it is a label we lay on, and the point carried its distinctness before any label touched it. The
separation and the name of the separation are not two things joined by an act of recording. They are one geometric fact seen
twice, and because they are one fact the label costs nothing intrinsic to keep.

To see why, ask what a coordinate actually \emph{is}. A chart $\varphi$ is a diffeomorphism from an open set of $M$ onto an
open set of $\mathbb{R}^n$, and $x^\mu(p) = \varphi(p)^\mu$ is simply the $\mu$-th slot of where $\varphi$ sends $p$. The
point $p$ is an element of the manifold with or without $\varphi$; the numbers $x^\mu(p)$ are the frame's reading of it, and
a different chart $\varphi'$ gives different numbers for the same point. This is the plain content of general covariance: the
labels are arbitrary and the point is not. When Einstein~\cite{einstein1916} built the field equations to hold in every
coordinate system, he was insisting exactly this --- that the coordinates carry no physics, that they are the tange we lay
on, contravariant and ours, and that the geometry underneath is indifferent to which labels we chose.

The consequence is that a difference requires no allocation of its own --- but the claim must be scoped exactly, because
scoped loosely it is false. The manifold does not set aside a store for the distinctness of each pair of points; the
distinctness is read straight off the smooth structure that was there anyway. It does not follow that the geometry carries
\emph{nothing}. It carries one standing datum before any measurement: the atlas itself, the fixed background of compatible
charts that lets a coordinate mean anything at all. That atlas is a single structural fixture, not a fresh cell spent on each
separation. So ``free'' here does not mean fast or cheap --- those would be claims about a clock, and this is not one. It
means \emph{nothing added}: the difference is decided where it stands, off structure already present, not recorded to be
retrieved. Difference, in this geometry, is free in the way a coordinate is free, and the only standing cost is the atlas
that was laid down once.

Underneath this is an idea older than the manifold: a symbol means nothing in isolation and means only by its contrast with
what it is not. Leibniz~\cite{leibniz1686} held that things sharing every property are one and the same --- the identity of
indiscernibles --- and the manifold takes him at his word, run backwards: two points are one until some open set tells them
apart, and a coordinate is worth exactly the contrast it draws. A label that discriminated nothing would carry no
information and name no point. The value of $x^\mu(p)$ is not the number itself but the fact that it differs from
$x^\mu(q)$; the coordinate is the contrast, not a tag pinned on after the contrast was noticed.

So take stock of what the geometry has handed you at the end of the ground floor: not a measurement, and not yet a number ---
a \emph{position}, which chart, which slot, relative to what. That is all of it, and it is exactly enough, because a position
told apart from other positions is the seed of everything the physics will grow. And the label is about to earn its keep a
second way: in the next chapter the geometry will \emph{funge} --- gather points that a coarse resolution cannot separate
into a single covariant reading --- and the size of that gathering will be its first number. The label is precisely what the
count will bag by. But that is the next chapter's work; here the only settled thing is that stepping onto the count carries
no hidden store to be paid for later. The single bill the instrument ever runs up is the turn, never the label.

\section{What a difference is not}

The ground floor is nearly built, and the last thing to do on it is the most easily skipped: to say what the first fact is
\emph{not}. A separation is not yet a distance. It is not a measurement, not a magnitude, not a count --- the smallest
possible geometric act, and no more. The temptation, having exhibited two distinct points, is to read more into them than
they hold; the whole discipline of this book is the refusal of that temptation, and the refusal is easiest to learn here,
where there is least to inflate.

The relation the geometry actually has is inequality, and inequality alone, and the cleanest way to see its poverty is to
list what a metric would add and note that none of it is present. A metric $g_{\mu\nu}$ supplies a length, $\mathrm{d}s^2 =
g_{\mu\nu}\,\mathrm{d}x^\mu\,\mathrm{d}x^\nu$; without it there is no $\mathrm{d}s$, no notion that $p$ and $q$ are near or
far, only that they are two. There are three refusals folded into this, and each rules out a structure the reader may be
tempted to assume is already there.

The first refusal: a difference is \emph{not an order}. To be a number is at the least to sit in a chain --- to have a
less-than and a greater-than. The bare manifold exposes no such thing. Distinctness is symmetric, $p \neq q$ iff $q \neq p$,
and blind to direction; it cannot say which point comes first, or lies higher, or is larger. Order is extra structure --- an
orientation, a time function, a sign convention --- laid on later, never a feature of the raw separation. A reader who
searches the ground floor for a which-is-bigger will not find one, and not because it is hidden: at this tier the question
does not so much as parse.

The second refusal: a difference is \emph{not a magnitude}. A measurement answers \emph{how much}; a separation answers only
\emph{whether}. There is no unit, no scale, no distance --- nothing that would let the geometry call two points far apart or
close, only that they are distinct. The datum is nominal in the plainest sense, same or different and no finer. Riemann, in
the lecture that founded the subject~\cite{riemann1854}, was careful to separate the bare manifold --- the
\emph{Mannigfaltigkeit}, the manifold of distinguishable positions --- from the metric relations one may afterward impose on
it; the magnitude belongs to the second, and the first fact lives entirely in the first.

The third refusal is the one this book is built to make precise, and it is cleanest stated in the two working verbs. A
difference is a pure \emph{tange} --- a contravariant selection, a label picking out \emph{this} against \emph{that}. A count
is something else: it is the size of a \emph{funge}, the covariant gathering of like things into one heap. You cannot take
the size of a heap before you have a heap, nor bag points by their sameness before you have first tanged them apart. So ``not
yet a number'' is not a shortfall the instrument will outgrow by trying harder; it is an ordering as strict as any in the
construction. Tange precedes funge; number lives wholly on the funge side of that line; and no amount of staring at a single
separation will yield a magnitude. That, in miniature, is the discipline of the whole book: hold exactly what is there --- a
distinctness, a position --- and decline to inflate it into a size it has not earned. The first fact keeps the same restraint
the final bracket will keep; between them stands only the long climb, and the climb begins when the geometry first funges its
tanged separations into a count.
